\section{Comparative Case Studies}
\label{sec:case-studies}
We select cases that expose distinct requirements of a specification test
oracle: arithmetic safety, quantified functional relations, and spatial
ownership. Each case identifies its source, input, observed tool outcome,
and control. These are recorded case studies rather than aggregate results
for the full benchmark. The evidence inventory pins the original sources,
logs, and proof receipts. The read cases use CN/Fulminate revision
\texttt{39efc873f} (2026-04-16); least-root and divisibility use
\texttt{181d04773} (2026-10-07). Both archived version strings carry a
\texttt{dirty} suffix, and binary hashes identify the tested builds.
These historical results remain separate from the latest-version evaluation
specified in Section~\ref{sec:evaluation}.

\subsection{An integer-boundary violation missed at runtime}
\label{sec:read-min}
The primary \texttt{read} program returns the value in an owned cell and
preserves that cell. Listing~\ref{lst:read-mutant} changes only the return
expression. The CN and QCP versions express the same ownership and value
requirements.
\begin{lstlisting}[caption={Negation mutant of read. The contract still requires the original value.},label={lst:read-mutant}]
int read(int *p)
/*@ With old
    Require data_at(p, int, old)
    Ensure __return == old &&
           data_at(p, int, old) */
{ return -*p; }
\end{lstlisting}
For a 32-bit signed integer, the input $\mathit{old}=-2^{31}$ makes the
mathematical negation $2^{31}$ unrepresentable. C classifies an expression
result outside its type's representable range as undefined behavior
(Section 6.5, paragraph 5 of the C11 draft~\cite{iso2011c}). The entry
cell is valid; it is the executed negation that violates safety.

The recorded Fulminate pipeline uses \texttt{cn instrument --run} with GCC.
Both the pipeline and a subsequent invocation of the generated executable
exit with code 0; the latter produces no diagnostic. \tespec{} reports a
satisfiable entry and an invalid safety obligation, yielding \fail{}.
For this input the range check specializes to the false proposition
\[
                    -2^{31}\le 2^{31}\le 2^{31}-1.
\]
The decisive evidence is the arithmetic-safety check, rather than a guessed
native return value. The archived log does not record that value.

\begin{table}[t]
\centering\small
\caption{Negation mutant: controls and the boundary discrepancy. Runtime rejection denotes Fulminate exit code 2.}
\label{tab:read-min}
\begin{tabular}{@{}lll@{}}
\toprule
Initial value & Fulminate & \tespec{}\\
\midrule
\texttt{INT\_MIN} & Accept (exit 0) & \fail{} (safety)\\
$0$ & Accept (exit 0) & \pass{}\\
$-1,1,7,\texttt{INT\_MAX}$ & Runtime rejection & \fail{}\\
\bottomrule
\end{tabular}
\end{table}
The controls in Table~\ref{tab:read-min} distinguish the missed safety
violation from ordinary incorrect return values. The correct read
implementation also passes at \texttt{INT\_MIN}. This is a missed violation
in the recorded runtime configuration. Compiler and sanitizer configuration
form part of that observation; a sanitizer-augmented baseline is a separate
configuration in the evaluation protocol. The \tespec{} result here is
automatic solver evidence, not an independently replayed Rocq certificate.

\subsection{A quantified minimum without a finite monitor}
For the \texttt{root\_cell} example, $n=13$ produces the value 4, while
Equation~\eqref{eq:least-root} still quantifies over all nonnegative integers.
The recorded CN instrumenter rejects the original postcondition because it
cannot infer an upper bound for the quantified variable. The same limitation
appears in the archived Bennet and Darcy configurations. Thus this case
exposes an assertion-support boundary, rather than an incorrect runtime
answer.

The \tespec{} postcondition-only replay initially returns \unknown{}.
Proof completion checks its complete group of 26 original obligations and
produces \pass{}. The proof reasons directly about the quantified relation.
As a control, the artifact proves its equivalence to the finite formula in
Equation~\eqref{eq:finite-root}; that formula passes Fulminate on the same
fixture. The comparison isolates the additional work required to obtain an
executable monitor from a valid mathematical specification.

The intermediate assertion provides a second observation site. Replacing
$\mathsf{LeastRoot}(n,r)$ by $\mathsf{LeastRoot}(n,r-1)$ preserves the program
execution but changes the assertion to $\mathsf{LeastRoot}(13,3)$. The archived
27-obligation group proves the other obligations and refutes this assertion
VC. Its receipt does not admit a complete execution-counterexample
certificate, so this is reported as a checked VC refutation. It demonstrates
that internal specification changes can be observed even when inputs and
outputs are unchanged.

\subsection{Existence and nonexistence in one contract}
The divisibility program returns whether \texttt{a \% b == 0}, under $b>0$.
Its mathematical specification distinguishes two proof tasks:
\[
(r=1\land\exists q\in\mathbb Z.\ a=qb)
\ \lor\ (r=0\land\forall q\in\mathbb Z.\ a\ne qb).
\]
The inputs $(6,3)$ and $(7,3)$ exercise both branches. The first proof supplies
$q=2$; the second excludes every integer witness. Both are positive
conformance tests. Automatic processing leaves both unresolved, and proof
completion checks each complete six-obligation execution group. The receipts
record their logical dependencies and the kernel-checking route.

The finite remainder-based CN contract passes both inputs. Adding its
logically redundant, unbounded universal consequence causes the tested
instrumenter to reject the assertion with a lower-bound inference diagnostic.
A separately kernel-checked Rocq lemma establishes the equivalence and the
redundancy of that consequence. Together with least-root, this case shows
that a formulation change can determine whether a runtime oracle is available,
even when the intended property is unchanged.

\subsection{A spatial mismatch resolved by either proof route}
Two read variants omit the owned cell from their return contract. The body
still retains that cell. In the exact resource interpretation used here,
a nonempty owned heap does not satisfy \texttt{emp}. An automatically added
frame for user heap would hide precisely this specification defect.

Across six integer inputs per variant, the recorded automatic run with
spatial SMT disabled leaves 12 cases \unknown{}. Independent proof-completion
jobs refute the original spatial obligations and check entry feasibility.
Enabling spatial SMT resolves those same 12 cases as \fail{}, with
countermodels checked against the finite-heap interpretation and agreement
with the existing kernel-checked refutations. The remaining 49 raw
classifications are unchanged, including all 22 passing cases. These solver
results retain their automatic-evidence label; agreement with a checked
refutation does not make the solver output itself a kernel certificate.

Fulminate also detects the missing-ownership variants. This control therefore
isolates the effect of \tespec{}'s proof backend: a residual spatial obligation
can represent a real defect, and a stronger solver can decide it without
changing execution or rejecting correct controls. The full read group contains
61 variant/input pairs across ten paired configurations. Fulminate matches
60 expected outcomes; \tespec{} with per-execution ITP or with spatial SMT
matches all 61 under the group's declared policy. Seven of these outcomes
are expected memory-access errors for specifications that supply no read
permission, and are kept separate from ordinary assertion failures.
